\thispagestyle{plain}
\chapter*{Resumo}

A crescente adoção de componentes de software de terceiros aumenta a superfície de ataque das cadeias de fornecimento, tornando difícil para as organizações compreender e mitigar os riscos introduzidos pelas suas dependências. O Software Bill of Materials (SBOM) emerge como um mecanismo de transparência, mas os formatos existentes carecem de ferramentas analíticas que transformem dados de inventário em inteligência de risco acionável.

Esta dissertação apresenta o SBOM Lens, uma framework baseada em grafos para avaliação de risco em cadeias de fornecimento de software. O sistema ingere SBOMs no formato CycloneDX e constrói um grafo de dependências orientado, sobre o qual aplica técnicas de ciência de redes para expor vulnerabilidades estruturais. O trabalho cobre quatro domínios principais: (1) a visualização interativa do grafo SBOM, que permite explorar e editar a topologia das dependências; (2) a gestão de vulnerabilidades, com integração em múltiplas fontes de dados como NVD, OSV e GitHub Security Advisories, incluindo o cálculo de scores CVSS v3.1; (3) a avaliação de risco estrutural, que identifica pontos de falha única, dependências circulares, conflitos de versão e concentração de dependências através de métricas de entropia e decomposição em núcleos; e (4) a análise de grafos, que aplica algoritmos de deteção de comunidades, decomposição k-core e análise espetral para modelar a propagação de vulnerabilidades.

O SBOM Lens demonstra que a teoria dos grafos e a ciência de redes constituem uma abordagem eficaz para priorizar riscos de segurança em cadeias de fornecimento de software modernas.

\keywordspt{SBOM, Segurança de Cadeias de Fornecimento de Software, Teoria dos Grafos, Análise de Vulnerabilidades, Avaliação de Risco, CycloneDX.}

\MediaOptionLogicBlank

\pdfbookmark[1]{Abstract}{abstract}
\chapter*{Abstract}

The widespread adoption of third-party software components continuously expands the attack surface of modern software supply chains, making it difficult for organisations to understand and mitigate the risks introduced by their dependencies. The Software Bill of Materials (SBOM) has emerged as a transparency mechanism, yet existing tooling largely treats SBOMs as flat inventories rather than structured graphs amenable to rigorous analysis.

This dissertation presents SBOM Lens, a graph-based framework for software supply chain risk assessment. The system ingests SBOMs in CycloneDX format and constructs a directed dependency graph, over which network science techniques are applied to expose structural vulnerabilities. The work covers four main domains: (1) interactive SBOM graph visualisation, enabling exploration and editing of the dependency topology; (2) vulnerability management, integrating multiple data sources — NVD, OSV, and GitHub Security Advisories — with full CVSS v3.1 base score computation; (3) structural risk assessment, which identifies single points of failure via articulation point analysis, circular dependencies, version conflicts through diamond dependency detection, and dependency concentration using entropy metrics and k-core decomposition; and (4) graph analysis, applying community detection, k-core decomposition, and spectral graph theory to model vulnerability propagation using epidemic spreading analogues.

SBOM Lens demonstrates that graph theory and network science provide an effective foundation for contextualising and prioritising security risk in modern software supply chains.

\keywordsen{SBOM, Software Supply Chain Security, Graph Theory, Vulnerability Analysis, Risk Assessment, CycloneDX.}

\MediaOptionLogicBlank